% A.70 reserved by the root.
\subsection[Finite noncongruence isolation of each realized point]{One explicit finite noncongruence neighbourhood isolates each realized point}

\begin{statement}{A.70}{Proved}
{Finite Kummer isolation in every fixed rational Legendre realization}
Let $\lambda_0=a/c\in\Q\setminus\{0,1\}$, where
$a+b=c$, $c>0$, and $a,b,c$ are pairwise coprime nonzero
integers.  Put $e_p=v_p(abc)$ and define
\[
 N_{\lambda_0}=
 \left(\prod_{\substack{p\mid abc\\p\ {
m odd}}}p^{e_p}\right)
 \begin{cases}
       2^{e_2-1},&e_2\ge5,\\
       1,&e_2\le4.
 \end{cases}
\tag{A.70.1}
\]
Write $N=N_{\lambda_0}$.
Among the noncuspidal Selmer sections with labeled integral
realization $E_{\lambda_0}$, the congruences
\[
             A_s\equiv\lambda_0,\qquad
             B_s\equiv1-\lambda_0
                       \pmod{(\widehat{\Q^\times})^N}
\tag{A.70.2}
\]
hold precisely for $s=s_{\lambda_0}$.
Equivalently, the explicit pointed finite etale cover
\[
 V_{\lambda_0,N}:\quad
       \lambda_0u^N+(1-\lambda_0)v^N=1,\quad uv\ne0,
 \qquad t=\lambda_0u^N,\qquad P=(1,1)
\tag{A.70.3}
\]
is a neighbourhood of exactly that section in this
realization class.

In particular, the point is isolated within its full
geometric congruence fibre by the one indicated finite
abelian enlargement.  More precisely, put
$E=\pi_1(U)$, $\Pi=\pi_1(U_{\overline\Q})$,
$T=\Pi^{\rm ab}$,
\[
 C_H=\ker(\Pi\to H),\qquad
 M_N=\ker(\Pi\to T/NT),\qquad
 E_{{\rm ab},N}=E/M_N.
\]
There is an actual topological isomorphism
\[
 \frac{E}{C_H\cap M_N}
       \simeq
 E_H\mathbin{\times}_{(H^{\rm ab}/NH^{\rm ab})\rtimes G_\Q}
                         E_{{\rm ab},N},
\tag{A.70.4}
\]
where the two maps to the common arithmetic quotient use
the point-$\lambda_0$ splitting.
The Selmer-section fibre over the image of $s_{\lambda_0}$
in (A.70.4) contains only $s_{\lambda_0}$.
This does not assert that the image in $E_H$ alone forces
the additional congruences (A.70.2).
\end{statement}
\begin{proof}
We first give the exact local arithmetic behind the
choice of $N$.  The Legendre formula in these integers is
\[
             j(E_{\lambda_0})=
             \frac{256(a^2-ac+c^2)^3}{a^2b^2c^2}.
\tag{A.70.5}
\]
At a prime dividing $abc$, exactly one of $a,b,c$ is
divisible by that prime, and the numerator factor
$a^2-ac+c^2$ is a unit.  Hence
\[
 v_p(j)=
 \begin{cases}
       -2e_p,&p\text{ odd and }p\mid abc,\\
       8-2e_2,&p=2.
 \end{cases}
\tag{A.70.6}
\]
At every other odd prime the $j$ invariant is integral.
Thus the set of potentially multiplicative primes is
exactly
\[
 S=\{p\mid abc:p\text{ odd}\}
                 \cup\{2:e_2\ge5\}.
\tag{A.70.7}
\]
At every prime outside $S$, A.37 excludes a locally
cuspidal restriction of any section in the stated
realization class.

For the three cusp labels define the actual rational
units on $U$
\[
        f_0(t)=1-t,\qquad f_1(t)=t,\qquad
                        f_\infty(t)=\frac{t}{t-1}.
\tag{A.70.8}
\]
Each $f_i$ is regular at its indicated cusp $i$ with
value one.  Its Kummer pullback along every section in
that cusp packet is therefore the trivial class; the
formal unit admits compatible roots with constant
term one.  At the realized point their values are
\[
                  f_0(\lambda_0)=b/c,\quad
                  f_1(\lambda_0)=a/c,\quad
                  f_\infty(\lambda_0)=-a/b.
\tag{A.70.9}
\]
At any $p\in S$, put $e=e_p$.
The following table lists their valuations and the
exact difference from one of the unit entry:
\[
\begin{array}{c|ccc|c}
 &v_p(f_0)&v_p(f_1)&v_p(f_\infty)&
                         \text{unit entry minus one}\\ \hline
 p\mid a&0&e&e&f_0-1=-a/c\\
 p\mid b&e&0&-e&f_1-1=-b/c\\
 p\mid c&-e&-e&0&f_\infty-1=-c/b.
\end{array}
\tag{A.70.10}
\]
In each row the expression in the last column has
valuation exactly $e$.

For odd $p\in S$ put $N_p=p^e$.
Neither entry of nonzero valuation in (A.70.10) is
an $N_p$th power, since $N_p>e$ and the valuation
of an $N_p$th power is divisible by $N_p$.
The remaining entry belongs to
$(1+p^e\Z_p)\setminus(1+p^{e+1}\Z_p)$.
If it were an $N_p$th power of a unit, the
root-of-unity component of that unit would be one:
raising to $p^e$ is an automorphism of $\mu_{p-1}$.
The exact principal-unit image is
\[
             (1+p\Z_p)^{p^e}=1+p^{e+1}\Z_p,
\tag{A.70.11}
\]
as follows from the logarithm isomorphism with
$p\Z_p$.  This also excludes the unit entry.
Thus all three values in (A.70.9) are nontrivial
in $\Q_p^\times/(\Q_p^\times)^{N_p}$.

For $p=2\in S$, one has $e\ge5$ and puts
$N_2=2^{e-1}$.
Again $N_2>e$ excludes the two entries of nonzero
valuation.  For units the decomposition
$\Z_2^\times=\{\pm1\}\times(1+4\Z_2)$ and
the logarithm on $1+4\Z_2$ give the exact image
\[
                  (\Z_2^\times)^{2^{e-1}}
                              =1+2^{e+1}\Z_2.
\tag{A.70.12}
\]
The unit entry in (A.70.10) has difference from one
of valuation exactly $e$, so it too is excluded.
This checks the dyadic case directly, with every
one of the three possible cusp labels retained.

Since $N_p\mid N$, the congruences (A.70.2)
imply equality of the local classes of all three
functions (A.70.8) with their point values modulo
$N_p$.  For $f_\infty$ the exact Kummer expression
is $-A_s/B_s$, so its fixed sign is retained.
If $s_p$ belonged to cusp packet $i$, the same
class would be one by regularity of $f_i$ there,
contradicting (A.70.10)--(A.70.12).
There are therefore no local cusps at any prime
in $S$, and A.37 excludes them outside $S$.
The no-cusp case of A.37 recognizes $s$ as a
rational point section.  The labeled integral
comparison identifies its parameter exactly with
$\lambda_0$ by the primitive-isogeny argument
at the end of A.54.
Conversely the point section has (A.70.2).
If $S$ is empty, (A.70.1) gives $N=1$ and the
same argument starts with no possible cuspidal
prime; A.37 and the labeled comparison apply
without any power test.  This includes the
zero-bad-prime case.

We check the stated cover and subgroup maps.
The two equations
$u^N=t/\lambda_0$, $v^N=(1-t)/(1-\lambda_0)$
make (A.70.3) finite etale of degree $N^2$ over
$U$, with geometric deck group $\mu_N^2$ and
the indicated rational point over $\lambda_0$.
The classes of $t$ and $1-t$ are independent
modulo $N$ over $\overline\Q$: their valuations
at zero and one respectively detect the two
exponents.  Thus this is the connected cover
whose geometric subgroup is $M_N$.
With the point-$\lambda_0$ splitting, its fibre
is the affine $G_\Q$-set $T/NT$.
A relative cocycle has a fixed sheet exactly
when its class in $H^1(G_\Q,T/NT)$ is zero.
These two relative Kummer classes are exactly
(A.70.2).  Choosing a fixed sheet identifies its
stabilizer with $\pi_1(V_{\lambda_0,N})$, which
proves the full-section lifting equivalence.

Finally both $C_H$ and $M_N$ are closed normal
geometric subgroups of $E$.  Their product is
closed, since it is the image of their compact
product under multiplication.
For any two such subgroups the actual quotient map
\[
 E/(C_H\cap M_N)\longrightarrow
          (E/C_H)\mathbin{\times}_{E/(C_HM_N)}(E/M_N)
\tag{A.70.13}
\]
is bijective.  Indeed, for representatives $e,f$
of a compatible pair write $e^{-1}f=cn$ with
$c\in C_H$, $n\in M_N$; the element $ec$
represents both required cosets.  Its kernel is
the indicated intersection.  Compactness and the
Hausdorff topology make this a topological
isomorphism.

The common quotient has geometric kernel
$H^{\rm ab}/NH^{\rm ab}$.
To verify this exact identification, A.56 gives
the surjection $T\to H^{\rm ab}$ and the
surjection $\Gamma_2\Pi\to[H,H]$.
Since $M_N$ is the inverse image of $NT$, its
image in $H$ contains $[H,H]$ and maps onto
$NH^{\rm ab}$; it is precisely the full inverse
image of this latter subgroup.
The point section splits the resulting arithmetic
extension, giving
$E/(C_HM_N)=(H^{\rm ab}/NH^{\rm ab})\rtimes G_\Q$.
Substituting into (A.70.13) proves (A.70.4).
Any Selmer section in its point fibre has both
the stated integral realization and (A.70.2),
so the preceding recognition proof applies.
A global cusp section cannot enter this fibre:
its Legendre representation has the invariant line
of its quadratic Tate filtration, whereas that of
the point $\lambda_0$ is absolutely irreducible by A.12.
\end{proof}

\paragraph{The exact finite kernel beyond the congruence image.}
Write $N=2^\alpha3^\beta M$ with $(M,6)=1$, and put
$x_+=\max(x,0)$.
The kernel of (A.70.4) over $E_H$ is the actual subgroup
\[
 K_N=\ker\bigl(T/NT\longrightarrow
                         H^{\rm ab}/NH^{\rm ab}\bigr).
\tag{A.70.14}
\]
A.56's character map
$(a,b)\mapsto(a\bmod8,b\bmod8,a+b\bmod3)$
gives the explicit decomposition
\[
 K_N\simeq
 (\Z/2^{(\alpha-3)_+})^2(1)
 \oplus\Z/3^\beta(1)
 \oplus\Z/3^{(\beta-1)_+}(1)
 \oplus(\Z/M)^2(1),
\tag{A.70.15}
\]
with trivial $\Z/1$ factors omitted.
At two the map is multiplication by eight on the
two coordinates when $\alpha>3$.
For $\beta\ge1$ the three-primary identification is
\[
 (u,v)\longmapsto(u,-u+3v)
 \quad\text{from}\quad
 \Z/3^\beta\oplus\Z/3^{\beta-1}
                   \quad\text{into }(\Z/3^\beta)^2.
\]
It is injective and has precisely the vectors whose
coordinate sum vanishes modulo three; its two
coordinates retain the cyclotomic action.
The remaining primary factors map identically to
$(\Z/M)^2$.  Thus the finite geometric degree of
this enlargement of $E_H$ is exactly
\[
                       |K_N|=
       2^{2(\alpha-3)_+}\,3^{(2\beta-1)_+}\,M^2.
\tag{A.70.16}
\]
For the particular modulus (A.70.1), the added
dyadic depth in each coordinate is $e_2-4$ when
$e_2\ge5$, and is zero otherwise.

\noindent\emph{The remaining congruence-recognition question.}
Equation (A.70.4) is a finite geometric enlargement
of $E_H$.  The theorem describes its point fibre
among Selmer sections and gives its specific
pointed cover.
For a section known only through its $E_H$ image,
the finite classes (A.70.2) are additional data.
A.56 calculates their common quotient with $H$,
and A.67 exhibits a local degree-five instance
which distinguishes two equal full congruence
images.  No implication supplying the missing
finite classes from $E_H$ alone is asserted.
